\input{estilo}
\usetikzlibrary{shapes.geometric}
% No está en la pizarra: los componentes no se corresponden con las máquinas.
\begin{document}
\begin{figura}\begin{tikzpicture}
  \draw[fill=white] (0,0) ellipse (3.8 and 2.35);
  % las máquinas, detrás
  \tikzset{maq/.style={draw=suave, dashed, fill=suave!12, rounded corners=3pt}}
  \draw[maq] (-2.95,-0.95) rectangle (-0.85,1.15);
  \draw[maq] (-0.45,0.2) rectangle (2.75,1.4);
  \draw[maq] (-0.45,-1.55) rectangle (2.75,-0.2);
  \node[etiq, text=suave, anchor=north west] at (-2.95,1.15) {máquina 1};
  \node[etiq, text=suave, anchor=north west] at (-0.45,1.4) {máquina 2};
  \node[etiq, text=suave, anchor=south west] at (-0.45,-1.55) {máquina 3};
  % los componentes
  \node[componente] (n1) at (-2.2,0.25) {};
  \node[componente] (n2) at (-1.45,-0.45) {};
  \node[componente] (n3) at (0.35,0.65) {};
  \node[componente] (n4) at (0.35,-0.75) {};
  % el componente que es otro sistema: una parte en cada máquina; sin relleno,
  % para que se vea el borde entre las dos máquinas
  \node[draw=tinta, ellipse, minimum width=12mm, minimum height=27mm] (big) at (1.95,-0.075) {};
  \node[componente] (b1) at (1.95,0.75) {};
  \node[componente] (b2) at (1.95,-0.9) {};
  \draw (b1) -- (b2);
  \draw (n1) -- (n2) (n1) -- (n3) (n2) -- (n4) (n3) -- (big) (n4) -- (big) (n3) -- (n4);
  % la interfaz
  \fill[tinta] ($(180:3.8 and 2.35)+(-0.11,-0.11)$) rectangle ++(0.22,0.22) coordinate[pos=0.5] (pa);
  \draw[suave] (pa) -- (n1);
  \draw[<->] ($(pa)+(-0.12,0)$) -- ++(-1.0,0);
  % anotaciones
  \node[etiq, align=right, anchor=east, text=suave] (t1) at (-3.6,2.0) {dos componentes\\[-1pt]en una máquina};
  \draw[suave] (t1.east) -- (-2.6,1.15);
  \node[etiq, align=left, anchor=west, text=suave] (t2) at (4.1,-0.5) {un componente que es\\[-1pt]otro sistema, repartido\\[-1pt]en dos máquinas};
  \draw[suave] (t2.west) -- (big.east);
\end{tikzpicture}\end{figura}
\end{document}
